% ==================== CHAPTER 3: MATERIALS AND METHODS ====================
\section{Materials and Methods}
\subsection{Overview of the Computational Pipeline}

The overall methodology of this proposed study will follow a systematic, five-phase computational pipeline, as illustrated in Figure~\ref{fig:workflow}. This workflow is designed to progress from raw biochemical data curation to the identification of novel, drug-like lead compounds through advanced graph-based deep learning.

\begin{figure}[htbp]
\centering
\begin{tikzpicture}[
    node distance=0.65cm,
    phase/.style={rectangle, rounded corners=6pt, draw=black, thick, fill=blue!20, minimum width=6.8cm, minimum height=0.72cm, align=center, font=\bfseries\small},
    process/.style={rectangle, rounded corners=4pt, draw=black, thick, fill=cyan!12, minimum width=6.3cm, minimum height=0.48cm, align=center, font=\small},
    final/.style={rectangle, rounded corners=6pt, draw=black, very thick, fill=green!25, minimum width=6.8cm, minimum height=0.75cm, align=center, font=\bfseries},
    arrow/.style={->, thick, >=latex, draw=black!80}
]

% Phase 1
\node[phase] (ph1) {Phase 1: Data Acquisition and Preprocessing};
\node[process, below=0.2cm of ph1] (p11) {ChEMBL Retrieval (CHEMBL203)};
\node[process, below=of p11] (p12) {pIC$_{50}$ Conversion \& Binary Labelling};

% Phase 2
\node[phase, below=0.6cm of p12] (ph2) {Phase 2: Molecular Graph Construction};
\node[process, below=0.2cm of ph2] (p21) {PyTorch Geometric Graphs};

% Phase 3
\node[phase, below=0.6cm of p21] (ph3) {Phase 3: Model Development \& Training};
\node[process, below=0.2cm of ph3] (p31) {Scaffold Split + AttentiveFP Training};

% Phase 4
\node[phase, below=0.6cm of p31] (ph4) {Phase 4: Evaluation \& Interpretability};
\node[process, below=0.2cm of ph4] (p41) {Metrics + GNNExplainer};

% Phase 5
\node[phase, below=0.6cm of p41] (ph5) {Phase 5: Virtual Screening};
\node[process, below=0.2cm of ph5] (p51) {Related Kinases Library (n=1,872)};
\node[process, below=of p51] (p52) {Inference \& Lead Prioritization};

% Final Output
\node[final, below=0.6cm of p52] (final) {Ranked EGFR Lead Candidates};

% Arrows
\draw[arrow] (ph1) -- (p11);
\draw[arrow] (p11) -- (p12);
\draw[arrow] (p12) -- (ph2);
\draw[arrow] (ph2) -- (p21);
\draw[arrow] (p21) -- (ph3);
\draw[arrow] (ph3) -- (p31);
\draw[arrow] (p31) -- (ph4);
\draw[arrow] (ph4) -- (p41);
\draw[arrow] (p41) -- (ph5);
\draw[arrow] (ph5) -- (p51);
\draw[arrow] (p51) -- (p52);
\draw[arrow] (p52) -- (final);

\end{tikzpicture}
\caption{Proposed five-phase computational pipeline for GNN-based EGFR inhibitor classification and virtual screening.}
\label{fig:workflow}
\end{figure}

\subsection{Data Acquisition and Preprocessing}

To build a robust predictive model, high-quality bioactivity data will be retrieved from the ChEMBL database (version 33 or later) via its RESTful API, specifically targeting the Epidermal Growth Factor Receptor (EGFR, ChEMBL ID: CHEMBL203). To ensure quantitative reliability, the dataset will be strictly filtered to include only records with exact Half-Maximal Inhibitory Concentration ($IC_{50}$) values, discarding any ambiguous measurements (e.g., `>', `<', or `$\sim$'). 

These values will be mathematically transformed into the negative logarithmic scale ($pIC_{50} = -\log_{10}(IC_{50} \times 10^{-9})$) to normalize the data distribution. A binary thresholding strategy will be implemented to frame the problem as a classification task: compounds with $pIC_{50} \geq 6.0$ (equivalent to $IC_{50} \leq 1 \mu M$) will be classified as active inhibitors, while those below this threshold will be deemed inactive.

Rigorous structural standardization will be performed using RDKit's \texttt{MolStandardize} module. This pipeline will strip disconnected fragments (salts and solvents), neutralize charges where appropriate, and generate canonical SMILES to identify and remove duplicated entries. In cases where multiple bioactivity values exist for a single canonical SMILES, the median $pIC_{50}$ will be retained. Finally, a physicochemical filter will isolate small molecules by restricting the molecular weight to a range of 100 to 1000 Daltons.

\subsection{Molecular Representation and Graph Construction}

Unlike traditional machine learning models that rely on fixed-length bit vectors (e.g., fingerprints), the proposed methodology will represent molecules natively as undirected graphs $\mathcal{G} = (\mathcal{V}, \mathcal{E})$, where $\mathcal{V}$ represents atoms (nodes) and $\mathcal{E}$ represents chemical bonds (edges). This transformation will be executed using PyTorch Geometric (PyG), enabling the preservation of complex topological and spatial relationships within the molecules.

\subsubsection{Node (Atom) Features}
Each node will be mathematically encoded as a comprehensive 29-dimensional continuous feature vector designed to capture fundamental biochemical properties essential for protein-ligand binding. The features will comprise:
\begin{itemize}
    \item \textbf{Atom type} (one-hot encoded, 10 dimensions): C, N, O, S, F, Cl, Br, I, P, and Unknown.
    \item \textbf{Degree} (one-hot encoded, 6 dimensions): Capturing the number of covalently bonded heavy-atom neighbours (0 to 5).
    \item \textbf{Formal charge} (1 dimension): A critical parameter for modelling electrostatic interactions within the EGFR kinase pocket.
    \item \textbf{Attached hydrogens} (one-hot encoded, 5 dimensions): Ranging from 0 to 4, acting as a proxy for hydrogen bond donor capacity.
    \item \textbf{Hybridization state} (one-hot encoded, 5 dimensions): SP, SP2, SP3, SP3D, SP3D2, which dictates the local 3D geometry of the atom.
    \item \textbf{Aromaticity} (binary, 1 dimension): Crucial for identifying potential $\pi-\pi$ stacking interactions with aromatic residues in the binding site.
    \item \textbf{Ring membership} (binary, 1 dimension).
\end{itemize}

\subsubsection{Edge (Bond) Features}
To further enrich the message-passing framework, each edge will be encoded as a 12-dimensional feature vector:
\begin{itemize}
    \item \textbf{Bond type} (one-hot encoded, 4 dimensions): Single, Double, Triple, Aromatic.
    \item \textbf{Conjugation} (binary, 1 dimension): Indicates extended electron delocalization.
    \item \textbf{Ring membership} (binary, 1 dimension).
    \item \textbf{Stereochemistry} (one-hot encoded, 6 dimensions): Captures E/Z and Cis/Trans configurations, which drastically affect ligand binding affinity.
\end{itemize}

\subsection{Data Splitting Strategy}

A major pitfall in cheminformatics machine learning is "scaffold memorization," where models perform exceptionally well on random test sets but fail to generalize to novel chemotypes. To rigorously evaluate the model's capacity for true virtual screening, a strict \textbf{Bemis-Murcko Scaffold Split} will be employed via RDKit. 

The Bemis-Murcko algorithm systematically strips away terminal side chains, leaving only the core ring systems and linker atoms (the scaffold). The dataset will be partitioned based on these unique scaffolds into an 80\% Training, 10\% Validation, and 10\% Test set. This guarantees that compounds sharing a core structural framework will be isolated into a single partition, forcing the model to evaluate entirely unseen molecular topologies during validation and testing.

\subsection{Model Architectures}

The core of this study will systematically evaluate four distinct Graph Neural Network (GNN) architectures functioning under the Message Passing Neural Network (MPNN) paradigm. The models to be developed include the Graph Convolutional Network (GCN), Graph Attention Network (GAT), Graph Isomorphism Network (GIN), and AttentiveFP.

\textbf{AttentiveFP} is proposed as the primary, focal architecture for this research. Unlike standard GCNs that aggregate neighbor information isotropically, AttentiveFP utilizes a multi-step, graph-level attention mechanism. Crucially, it provides native support for integrating multidimensional edge features directly into the message-passing update equations. This capability is highly advantageous for molecular property prediction, as the specific nature of a chemical bond profoundly influences the electronic and spatial profile of a pharmacophore.

To establish a baseline and demonstrate the superiority of graph-based methods over established cheminformatics techniques, Random Forest (RF) and Logistic Regression models will also be trained. These baselines will utilize traditional 2048-bit ECFP4 (Morgan) fingerprints with a radius of 2 bonds, generated via RDKit.

\subsection{Training and Hyperparameter Optimization}

All GNN models will be trained using PyTorch. To address the anticipated class imbalance in the ChEMBL dataset (approximately 74\% active compounds), a Weighted Binary Cross-Entropy (BCE) loss function with logits will be formulated. A positive class weight ($pos\_weight = 0.35$) will be applied to penalize misclassifications of the minority (inactive) class more heavily, ensuring the model does not bias its predictions toward the majority class.

The network weights will be optimized using the Adam optimizer. To ensure stable convergence, a \texttt{ReduceLROnPlateau} learning rate scheduler will be implemented (reduction factor = 0.5, patience = 10 epochs). An early stopping mechanism based on the Validation AUROC (patience = 20-30 epochs) will be utilized to prevent overfitting to the training scaffolds.

Hyperparameter optimization will be executed using the \textbf{Optuna} framework, leveraging the Tree-structured Parzen Estimator (TPE) algorithm for Bayesian optimization. The search space for the AttentiveFP model will be highly constrained, iterating over hidden dimensions, the number of GNN message-passing layers, the number of attention timesteps, dropout rates, batch sizes, learning rates, and L2 weight decay coefficients. The configuration yielding the highest AUROC on the validation set will be selected as the final model.

\subsection{Evaluation Metrics}

Due to the inherent class imbalance, relying solely on classification accuracy is insufficient and potentially misleading. Before defining the metrics, we establish the four fundamental outcomes of binary classification: true positives ($TP$) and true negatives ($TN$) denote correctly classified active and inactive compounds, respectively, while false positives ($FP$) and false negatives ($FN$) denote the corresponding misclassifications. From these quantities we derive the following intermediate measures:
\begin{align}
    \text{Precision} &= \frac{TP}{TP + FP}, \\
    \text{Recall} = \text{TPR} &= \frac{TP}{TP + FN}, \\
    \text{FPR} &= \frac{FP}{FP + TN},
\end{align}
where TPR and FPR denote the true positive rate (sensitivity) and false positive rate, respectively.

Model performance will be assessed primarily using threshold-independent metrics:
\begin{itemize}
    \item \textbf{Area Under the Receiver Operating Characteristic Curve (AUROC):} Measures the model's ability to distinguish between active and inactive classes across all possible classification thresholds. The ROC curve plots TPR against FPR as the decision threshold varies, and the AUROC is the area beneath it:
    \begin{equation}
        \text{AUROC} = \int_{0}^{1} \text{TPR}\,(\text{FPR}^{-1}(x))\; dx.
    \end{equation}
    Equivalently, the AUROC equals the probability that a randomly chosen active compound is ranked above a randomly chosen inactive one:
    \begin{equation}
        \text{AUROC} = \Pr\!\left(\hat{s}^{+} > \hat{s}^{-}\right),
    \end{equation}
    where $\hat{s}^{+}$ and $\hat{s}^{-}$ are the predicted scores for a positive and a negative instance, respectively.

    \item \textbf{Area Under the Precision-Recall Curve (AUPRC):} Heavily prioritizes the model's predictive precision specifically regarding the minority class, making it the most stringent metric for imbalanced data. It summarizes the trade-off between precision and recall across thresholds and is computed in practice via the average-precision estimator:
    \begin{equation}
        \text{AUPRC} = \sum_{n} \left( R_{n} - R_{n-1} \right) P_{n},
    \end{equation}
    where $P_{n}$ and $R_{n}$ denote the precision and recall at the $n$-th threshold.
\end{itemize}

Secondary metrics, evaluated at a default probability threshold of $0.5$, will include the F1-Score, the Matthews Correlation Coefficient (MCC)—which provides a balanced measure of true and false positives/negatives—and overall Accuracy. The F1-Score is the harmonic mean of precision and recall:
\begin{equation}
    \text{F1} = 2 \cdot \frac{\text{Precision} \cdot \text{Recall}}{\text{Precision} + \text{Recall}} = \frac{2\,TP}{2\,TP + FP + FN}.
\end{equation}
The MCC returns a value in $[-1, 1]$, where $1$ indicates perfect prediction, $0$ no better than random, and $-1$ total disagreement:
\begin{equation}
    \text{MCC} = \frac{TP \cdot TN - FP \cdot FN}{\sqrt{(TP + FP)(TP + FN)(TN + FP)(TN + FN)}}.
\end{equation}
Finally, overall Accuracy is the proportion of correctly classified compounds:
\begin{equation}
    \text{Accuracy} = \frac{TP + TN}{TP + TN + FP + FN}.
\end{equation}
All final reported metrics will be calculated exclusively on the held-out Bemis-Murcko test set.

\subsection{Model Interpretability}

To overcome the "black-box" criticism frequently leveled at deep learning in drug discovery, \textbf{GNNExplainer} will be applied to the trained AttentiveFP model. This post-hoc explainability algorithm operates by maximizing the mutual information between the model's prediction and a compact subgraph structure.

GNNExplainer will generate node and edge importance masks, visually highlighting the specific atomic neighborhoods most influential in driving a positive prediction. Explanations will be deliberately generated for widely studied, clinically approved EGFR inhibitors (e.g., Gefitinib, Erlotinib, Osimertinib, and Lapatinib). The resulting computational attributions will be critically analyzed against established EGFR kinase pharmacophore models (verifying if the model successfully learned the importance of hinge-binding quinazoline nitrogens or hydrophobic halogen substituents). This step will validate that the model is learning genuine biochemical principles rather than spurious dataset correlations.

\subsection{Virtual Screening Pipeline}

The culmination of the study will involve deploying the optimized and validated AttentiveFP model in a prospective virtual screening pipeline. A specialized screening library will be curated by retrieving compounds reported as active against phylogenetically and structurally related kinases (e.g., VEGFR2, BRAF, CDK2, ABL1, and MET) from ChEMBL. Any compounds overlapping with the primary EGFR training set will be strictly excluded to guarantee external, zero-shot validation.

Prior to inference, the screening library will be filtered to ensure baseline drug-likeness. Only compounds with a molecular weight between 250 and 600 Da, and a calculated LogP between -1 and 7, will be retained. The AttentiveFP model will then perform forward-pass inference on this unseen library (estimated $n=1,872$), generating a predicted probability of EGFR inhibition for each compound.

To triage the highest-scoring candidates for potential future experimental validation, a suite of ADMET-proxy properties will be calculated using RDKit. These will include Molecular Weight (MW), calculated partition coefficient (LogP), the number of hydrogen bond donors and acceptors, Topological Polar Surface Area (TPSA), strict compliance with Lipinski’s Rule of Five, and the Quantitative Estimate of Drug-likeness (QED) score. Candidates exhibiting both a high predictive probability ($>0.85$) and an optimal QED score will be proposed as high-value, novel lead compounds.

\subsection{Software and Implementation Details}

All computational modeling and data processing will be performed using open-source Python (version 3.10+) libraries to ensure complete reproducibility. Core packages will include RDKit (2023.09+) for molecular standardization and descriptor generation, PyTorch and PyTorch Geometric for GNN construction and training, Optuna for Bayesian hyperparameter tuning, and scikit-learn for baseline modeling and metric computation. Visualizations, including GNNExplainer outputs, will be generated using Matplotlib and Seaborn. All deep learning experiments will be executed on GPU-accelerated hardware (NVIDIA CUDA architecture) to facilitate the computationally intensive training of graph neural networks.

\newpage